\SourceLeaf{239}%
atque perpendimus, valores differentialium \(\mathrm{d}t,\mathrm{d}u\) sic prodeuntium. aequales esse debere, independenter a differentialibus \(\mathrm{d}x,\mathrm{d}y\), quantitatibus \(T\mathrm{d}x+U\mathrm{d}y\), \(U\mathrm{d}x+V\mathrm{d}y\) resp. inveniemus, post quasdam transformationes satis obvias:
\[
\begin{aligned}
C^3T&=\alpha Ab'b'+\beta Bb'b'+\gamma Cb'b'\\
&\quad-2\alpha'Abb'-2\beta'Bbb'-2\gamma'Cbb'\\
&\quad+\alpha''Abb+\beta''Bbb+\gamma''Cbb\\
C^3U&=-\alpha Aa'b'-\beta Ba'b'-\gamma Ca'b'\\
&\quad+\alpha'A(ab'+ba')+\beta'B(ab'+ba')+\gamma'C(ab'+ba')\\
&\quad-\alpha''Aab-\beta''Bab-\gamma''Cab\\
C^3V&=\alpha Aa'a'+\beta Ba'a'+\gamma Ca'a'\\
&\quad-2\alpha'Aaa'-2\beta'Baa'-2\gamma'Caa'\\
&\quad+\alpha''Aaa+\beta''Baa+\gamma''Caa
\end{aligned}
\]
Si itaque brevitatis caussa statuimus
\[
\begin{aligned}
A\alpha+B\beta+C\gamma&=D\qquad .\ .\ .\ .\ .\ .\ .\ .\ .\ .\quad(1)\\
A\alpha'+B\beta'+C\gamma'&=D'\qquad .\ .\ .\ .\ .\ .\ .\ .\ .\ .\quad(2)\\
A\alpha''+B\beta''+C\gamma''&=D''\qquad .\ .\ .\ .\ .\ .\ .\ .\ .\ .\quad(3)
\end{aligned}
\]
fit
\[
\begin{aligned}
C^3T&=Db'b'-2D'bb'+D''bb\\
C^3U&=-Da'b'+D'(ab'+ba')-D''ab\\
C^3V&=Da'a'-2D'aa'+D''aa
\end{aligned}
\]
Hinc invenimus, evolutione facta,
\[
C^6(TV-UU)=(DD''-D'D')(ab'-ba')^2=(DD''-D'D')CC
\]
et proin formulam pro mensura curvaturae
\[
k=\frac{DD''-D'D'}{(AA+BB+CC)^2}
\]
\Article{11}
Formulae modo inventae iam aliam superstruemus, quae inter fertilissima theoremata in doctrina de superficiebus curvis referenda est. Introducamus sequentes notationes:
